Yes. The distinguished extension makes a multiscale prefix compression work, and it gives a genuine constant-factor entropy reduction.

Assume a counterexample gives a coherent distinguished extension
[
L^\star=(v_1,\dots,v_n)
]
such that for every (i<j),
[
\Pr_L[v_j<_L v_i]\le \frac13.
\tag{1}
]
For comparable pairs the probability is (0); for incomparable pairs this is exactly the bias alignment assumption.

For clarity take (n=2^r).

### Dyadic prefix compression

Put the elements in their (L^\star)-order and recursively bisect:

[
[1,n]\to [1,n/2]\cup[n/2+1,n]\to\cdots.
]

At every node (B=A\sqcup C), restrict (L) to the elements of (B), and record only the binary merge word
[
W_B\in{A,C}^{|B|},
]
telling us whether each successive element came from the left or right half.

Doing this recursively is lossless:
[
\boxed{L\longleftrightarrow (W_B)*{B\in T*{\rm int}}.}
]

This is just the information recorded by merge sort.

Now comes the useful part.

For a balanced node (|A|=|C|=m), define
[
K_B=\operatorname{inv}(W_B)
=#{(a,c)\in A\times C:c<_L a}.
]

Every pair (v_i,v_j) belongs to opposite halves at exactly one node—its least common ancestor. Therefore

[
\boxed{
\operatorname{inv}_{L^\star}(L)=\sum_B K_B.
}
\tag{2}
]

So inversion distance has been decomposed **exactly by scale**.

Even better, (K_B) is literally a prefix statistic. If (d_t(W_B)) is the number of elements on the wrong side of the (t)-th prefix relative to the canonical word
[
A^mC^m,
]
then
[
\boxed{K_B=\sum_{t=1}^{2m-1}d_t(W_B).}
\tag{3}
]

It is the area between the random prefix path and the distinguished-prefix path.

---

### The bias gives entropy loss at every scale

From (1),
[
\mathbb E K_B
=============

\sum_{a\in A,c\in C}\Pr(c<_La)
\le \frac{m^2}{3}.
\tag{4}
]

For comparison, an unconstrained uniformly random interleaving has
[
\mathbb E K_B=\frac{m^2}{2}.
]

So at **every node of the hierarchy**, the merge path is macroscopically displaced toward (L^\star).

Here is an elementary entropy lemma.

If (W) is any random binary word containing (m) (A)'s and (m) (C)'s and
[
\mathbb E\operatorname{inv}(W)\le m^2/3,
]
then
[
\boxed{
H(W)
\le
2m-\frac{m^3}{3\ln2,(4m^2-1)}
\le
\left(1-\frac1{24\ln2}\right)2m.
}
\tag{5}
]

Numerically,
[
\frac1{24\ln2}=0.06011\ldots,
]
so
[
\boxed{H(W)\le0.9399,|W|.}
]

The proof is short. Let (X_t) indicate that position (t) contains an (A), and (p_t=\Pr(X_t=1)). Then
[
H(W)\le\sum_t h_2(p_t).
]
The inversion bound says the centre of mass of the (A)'s is shifted left by at least (m^2/6). Cauchy–Schwarz forces
[
\sum_t(p_t-\tfrac12)^2=\Omega(m),
]
and binary Pinsker,
[
1-h_2(p)\ge\frac{2}{\ln2}(p-\tfrac12)^2,
]
gives (5).

The optimal logistic calculation improves (0.9399) slightly to about (0.9383), but the elementary constant is enough.

---

### Sum over scales

At each dyadic level, the node sizes sum to (n). There are (\log_2n) levels. Therefore

[
\sum_B |B|=n\log_2n.
]

For uniform (L\in\mathcal L(P)),
[
H(L)=\log_2 e(P).
]

Subadditivity plus (5) gives

[
\boxed{
\log_2e(P)
\le
\left(1-\frac1{24\ln2}\right)n\log_2n.
}
\tag{6}
]

Thus approximately

[
\boxed{
\log_2e(P)\le0.9399,n\log_2n.
}
]

Equivalently,
[
\boxed{
e(P)\le (n!)^{,0.9399+o(1)}.
}
\tag{7}
]

Balanced non-dyadic trees give the same asymptotic statement.

So **yes: the coherent (2/3) orientation alone gives a constant (\epsilon<1) lossless entropy compression**, on the natural (n\log n) scale.

This does *not* follow merely from
[
\mathbb E\operatorname{inv}(L,L^\star)\le\frac13\binom n2.
]
A single global inversion bound still permits almost (n\log n) bits. What matters is that the (1/3) constraint holds for **every pair**, hence every dyadic block at every scale gets the (1/3) area bound. We lose (\Theta(n)) bits at each of (\Theta(\log n)) scales.

That is the key phenomenon.

### And it preserves BK structure unusually well

There is another exact property.

Take a BK edge swapping adjacent (x,y). There is a unique node (B) of the dyadic tree whose two children separate (x) and (y). Under the merge representation:

* every other (W_{B'}) stays unchanged;
* (W_B) changes by one adjacent
  [
  AC\leftrightarrow CA.
  ]

Thus

[
\boxed{\text{one BK edge }\longmapsto\text{ one local move at one scale}.}
\tag{8}
]

Each individual binary merge space is the lattice of monotone paths in an (m\times m) rectangle, equivalently an order-ideal lattice; covering graphs of distributive lattices are median graphs. ([arXiv][1])

So this compression gives almost exactly the bridge you were asking for:

[
\boxed{
\begin{array}{ccc}
\text{linear extension}
&\longrightarrow&
\text{hierarchy of median/cubical path spaces}[2mm]
\text{BK swap}
&\longrightarrow&
\text{one local edge at one scale}\
\text{inversion statistic}
&\longrightarrow&
\text{sum of prefix areas}\
2/3\text{ bias}
&\longrightarrow&
\text{constant entropy loss at every scale}.
\end{array}}
]

The general connection between entropy and the number of linear extensions is of course classical in sorting under partial information; (H=\log e(P)) is the information-theoretic quantity one expects to compress. ([arXiv][2]) Binary-word inversion numbers are also naturally encoded by Gaussian binomial coefficients. ([arXiv][3])

The part I find most promising is (8): **we have not merely compressed the vertices. We have produced a multiscale representation in which BK edges localize to one median factor.** The next question should be whether the Dirichlet form of the standard statistic admits a corresponding scale decomposition. If it does, the (6%) entropy saving itself is secondary—the real gain is that the spectral problem has been converted into a multiscale family of median-graph problems.

[1]: https://arxiv.org/pdf/0705.1025?utm_source=chatgpt.com "Recognizing Partial Cubes in Quadratic Time"
[2]: https://arxiv.org/abs/0911.0086?utm_source=chatgpt.com "Sorting under Partial Information (without the Ellipsoid Algorithm)"
[3]: https://arxiv.org/pdf/2605.25515?utm_source=chatgpt.com "Lipschitz Functions on Sparse Graphs II"
